\documentclass[11pt,a4paper]{beamer}
\usetheme{Madrid}
\usepackage[utf8]{inputenc}
\usepackage{amsmath}
\usepackage{amsfonts}
\usepackage{amssymb}
\usepackage{graphicx}
\usepackage{xcolor}
\usepackage{tikz}
\usepackage{booktabs}
\usepackage{array}
\usepackage{colortbl}
\usepackage{multicol}

% ============================================================
% TINY FONT FOR MAXIMUM CONTENT
% ============================================================
\usefonttheme{professionalfonts}
\setbeamerfont{normal text}{size=\scriptsize}
\setbeamerfont{itemize/enumerate body}{size=\scriptsize}
\setbeamerfont{itemize/enumerate subbody}{size=\tiny}
\setbeamerfont{block body}{size=\scriptsize}
\setbeamerfont{block title}{size=\scriptsize}
\setbeamerfont{frametitle}{size=\small}
\setbeamerfont{title}{size=\normalsize}
\setbeamerfont{subtitle}{size=\small}
\setbeamerfont{author}{size=\tiny}
\setbeamerfont{date}{size=\tiny}

\renewcommand{\scriptsize}{\fontsize{6pt}{7pt}\selectfont}
\renewcommand{\footnotesize}{\fontsize{6.5pt}{8pt}\selectfont}

\newcommand{\redflag}[1]{\textcolor{red}{\textbf{[!] #1}}}
\newcommand{\greennote}[1]{\textcolor{green!50!black}{\textbf{[Right] #1}}}
\newcommand{\examfocus}[1]{\textcolor{orange!80!black}{\textbf{[FOCUS] #1}}}
\newcommand{\trap}[1]{\textcolor{purple}{\textbf{[TRAP] #1}}}
\newcommand{\correct}[1]{\textcolor{green!50!black}{\textbf{[CORRECT] #1}}}
\newcommand{\scenario}[1]{\textcolor{blue!70!black}{\textbf{[SCENARIO] #1}}}

\title[CAMS Domain A - Missed Hard Questions]{CAMS Exam: The Missed Hard Questions\\High-Difficulty Concepts from Recent Test-Taker Feedback}
\author{Based on Post-Exam Reports \& Reddit Feedback}
\date{}

\setbeamercovered{transparent}
\setbeamertemplate{navigation symbols}{}
\setbeamertemplate{frametitle continuation}{}
\setlength{\parskip}{2pt}
\setlength{\itemsep}{1pt}

\begin{document}

\begin{frame}
\titlepage
\centering
\tiny Focus: FATF Immediate Outcomes (IOs), Insurance ML, Patriot Act Sections, DeFi Control, Nested Correspondent, and the "Most Correct" Answer Trap
\end{frame}

% ============================================================
% SECTION 1: FATF IMMEDIATE OUTCOMES (MISSED IN BOTH FILES)
% ============================================================

\begin{frame}{FATF Immediate Outcomes (IOs) - MISSING in your slides}
\begin{block}{What are Immediate Outcomes?}
FATF measures effectiveness through 11 IOs, NOT just technical compliance. The exam tests scenario-to-IO matching.
\end{block}
\begin{tabular}{p{0.35\textwidth}|p{0.60\textwidth}}
\hline
\rowcolor{lightgray}
\textbf{IO Number \& Name} & \textbf{What It Measures / Exam Scenario Example} \\
\hline
\textbf{IO.1: ML/TF Risks} & Country understands its risks (e.g., completed NRA) \\
\hline
\textbf{IO.2: International Cooperation} & Exchanging info with foreign FIUs; extradition success \\
\hline
\textbf{IO.3: Supervision} & FIs supervised and penalized for non-compliance \\
\hline
\textbf{IO.4: Preventive Measures} & FIs apply CDD, SARs, record-keeping effectively \\
\hline
\textbf{IO.5: Legal Persons \& Arrangements} & BO information available and accurate (lack of registry = failure here) \\
\hline
\textbf{IO.6: Financial Intelligence} & FIU produces and disseminates actionable intelligence \\
\hline
\textbf{IO.7: ML Investigation \& Prosecution} & ML prosecuted; proceeds confiscated \\
\hline
\textbf{IO.8: TF Investigation \& Prosecution} & TF prosecuted; no minimum threshold \\
\hline
\textbf{IO.9: Financial Sanctions} & Assets frozen without delay; reporting enforced \\
\hline
\textbf{IO.10: Confiscation} & Proceeds and instrumentalities confiscated \\
\hline
\textbf{IO.11: NPO Vulnerabilities} & Targeted supervision without disrupting legitimate NPOs \\
\hline
\end{tabular}
\redflag{Exam QUESTION: "Which IO is failing if no BO registry exists?" → IO.5}
\end{frame}

\begin{frame}{FATF IO Scenario Application}
\scenario{Country A has excellent ML/TF laws on paper. However, the FIU never shares reports with prosecutors. LE cannot obtain bank records without a 6-month court order. No foreign MLAT requests have ever been answered.}
\vspace{0.2cm}
\textbf{Q: Which FATF Immediate Outcome is MOST deficient?}
\begin{itemize}
\item A) IO.1 (Risk Understanding)
\item B) IO.4 (Preventive Measures)
\item C) IO.2 (International Cooperation)
\item D) IO.6 (Financial Intelligence)
\end{itemize}
\trap{Distractor: IO.4 seems correct because "preventive measures" includes SARs. But SARs exist; they just aren't used.}
\correct{Correct Answer: D (IO.6 - Financial Intelligence) because intelligence is not disseminated to LE. Also C (IO.2) is partially correct, but IO.6 is the primary failure.}
\examfocus{Exam often asks for the MOST deficient IO in a scenario. Know all 11 by name and core function.}
\end{frame}

% ============================================================
% SECTION 2: LIFE INSURANCE (UNDER-EMPHASIZED)
% ============================================================

\begin{frame}{Life Insurance - CDD Triggers \& ML Methods}
\begin{block}{CDD Triggers (FATF Rec. 12)}
CDD applies to life insurance policies when:
\begin{itemize}
\item Annual premium exceeds USD/EUR 10,000
\item Single premium exceeds USD/EUR 10,000
\item ANY suspicious transaction regardless of amount
\end{itemize}
\redflag{BENEFICIARY CDD: Must identify beneficial owner of a life insurance policy ONCE THE BENEFICIARY IS DESIGNATED (not necessarily at purchase).}
\end{block}
\begin{block}{Money Laundering via Life Insurance}
\textbf{Method 1 - Policy Loan Layering:}
\begin{enumerate}
\item Purchase policy with illicit funds (single premium)
\item Take a policy loan for 80-90\% of cash value
\item Loan check is "legitimate" - layered funds
\end{enumerate}
\textbf{Method 2 - Early Surrender:}
\begin{enumerate}
\item Purchase policy with dirty money
\item Surrender policy after short period (e.g., 6 months)
\item Receive surrender check minus penalty (the "laundering fee")
\end{enumerate}
\end{block}
\end{frame}

\begin{frame}{Life Insurance Exam Scenario - Policy Loan}
\scenario{A foreign PEP purchases a \$2.5 million single-premium life insurance policy. Three months later, he takes a \$2.1 million policy loan. The loan proceeds are wired to a real estate closing attorney for a luxury condo purchase. The PEP surrenders the policy one year later, receiving the remaining cash value.}
\vspace{0.2cm}
\textbf{Q: Which stage of money laundering is the policy loan?}
\begin{itemize}
\item A) Placement
\item B) Layering
\item C) Integration
\item D) Structuring
\end{itemize}
\trap{The loan is from a legitimate insurance company, so it appears clean. But the underlying premium was illicit.}
\correct{Correct Answer: B (Layering). The loan converts the illicit premium into a "legitimate" check from an insurance company, obscuring the origin. The surrender is integration.}
\redflag{The insurance company must apply EDD to the foreign PEP policyholder AND identify the beneficiary (even if an entity).}
\end{frame}

% ============================================================
% SECTION 3: USA PATRIOT ACT SECTIONS (MISSING DETAIL)
% ============================================================

\begin{frame}{USA PATRIOT Act - Key Sections Tested}
Test-takers report multiple questions on specific Patriot Act sections. Memorize these:
\begin{tabular}{p{0.20\textwidth}|p{0.75\textwidth}}
\hline
\rowcolor{lightgray}
\textbf{Section} & \textbf{Key Provision} \\
\hline
\textbf{Section 311} & Designate jurisdictions, institutions, or transactions as "primary money laundering concern" → special measures \\
\hline
\textbf{Section 312} & Enhanced Due Diligence for correspondent accounts with foreign banks (especially those with shell banks or PEPs) \\
\hline
\textbf{Section 313} & \textbf{PROHIBITS} U.S. banks from maintaining correspondent accounts for foreign shell banks (no physical presence, not affiliated with regulated bank) \\
\hline
\textbf{Section 314(a)} & LE can request information from FIs; FIs must search records and respond \\
\hline
\textbf{Section 314(b)} | \textbf{ALLOWS information sharing between FIs} (safe harbor from liability) \\
\hline
\textbf{Section 319(b)} | Requires respondent banks to identify BOs; U.S. courts can order forfeiture of funds in correspondent accounts \\
\hline
\textbf{Section 326} & Customer Identification Program (CIP) requirements \\
\hline
\textbf{Section 352} & Mandatory AML compliance programs \\
\hline
\textbf{Section 5318(h)} & Due diligence for foreign correspondent accounts \\
\hline
\end{tabular}
\redflag{Sections 313 (shell bank prohibition) and 314(b) (FI-to-FI sharing) are frequently tested.}
\end{frame}

\begin{frame}{Patriot Act Section 314(b) - Exam Scenario}
\scenario{Bank A suspects a customer is using multiple accounts across Bank B, Bank C, and Bank D to layer funds. Bank A wants to ask the other banks if they have seen similar activity.}
\vspace{0.2cm}
\textbf{Q: Under the USA Patriot Act, what allows Bank A to share information with Bank B?}
\begin{itemize}
\item A) Section 311
\item B) Section 313
\item C) Section 314(a)
\item D) Section 314(b)
\end{itemize}
\trap{Section 314(a) is for law enforcement requests, not voluntary FI sharing.}
\correct{Correct Answer: D (Section 314(b)). This section provides a safe harbor for FIs to share information regarding suspected ML/TF. Both parties must file SARs.}
\examfocus{314(a) = LE to FI. 314(b) = FI to FI.}
\end{frame}

% ============================================================
% SECTION 4: DeFi AND VASP CONTROL (ADVANCED)
% ============================================================

\begin{frame}{DeFi Protocol - Who is a VASP? (FATF Guidance)}
\begin{block}{The "Control or Influence" Test}
A person or entity that maintains \textbf{control or sufficient influence} over a DeFi protocol may be considered a VASP subject to AML obligations.
\vspace{0.2cm}
\textbf{Indicators of Control:}
\begin{itemize}
\item Owns >50\% of governance tokens (can unilaterally pass proposals)
\item Ability to upgrade smart contract code
\item Controls admin keys or multi-signature keys
\item Receives fees or revenue from protocol operations
\item Markets the protocol to users
\end{itemize}
\end{block}
\scenario{A software developer creates a DEX. He retains 68\% of governance tokens and holds the only admin key that can pause the contract. He does not collect fees directly.}
\vspace{0.1cm}
\textbf{Q: Is he a VASP under FATF guidance?}
\trap{He doesn't collect fees - but that's irrelevant.}
\correct{Correct Answer: YES. He has control via tokens and admin key. He must register and implement AML/CFT controls.}
\end{frame}

% ============================================================
% SECTION 5: NESTED CORRESPONDENT (EXPANDED)
% ============================================================

\begin{frame}{Nested Correspondent Accounts - The Complete Risk}
\begin{block}{Definition}
A "nested" account is when a respondent bank allows \textbf{its own downstream correspondent banks or third-party clients} to route transactions through the respondent's account at the correspondent bank, WITHOUT the correspondent's knowledge or consent.
\end{block}
\begin{tabular}{p{0.30\textwidth}|p{0.65\textwidth}}
\hline
\rowcolor{lightgray}
\textbf{Risk} & \textbf{Why It Matters} \\
\hline
Loss of visibility & Correspondent cannot see true originator/beneficiary \\
\hline
Sanctions evasion & Downstream party may be sanctioned entity \\
\hline
PEP exposure & Downstream party may be PEP with no EDD \\
\hline
TBML & Nested structure perfect for trade-based layering \\
\hline
Regulatory violation & Correspondent's CDD is effectively useless \\
\hline
\end{tabular}
\vspace{0.2cm}
\redflag{If a correspondent discovers undisclosed nesting, it must file a SAR, send an RFI, and consider termination. The respondent bank may be violating its own AML program.}
\end{frame}

% ============================================================
% SECTION 6: "MOST CORRECT" ANSWER TRAP
% ============================================================

\begin{frame}{The "Most Correct" Answer - CAMS Question Structure}
\begin{block}{Why This is the 1 Complaint}
Many CAMS questions present \textbf{two or three factually correct answers}. You must select the \textbf{BEST} or \textbf{MOST} correct based on FATF standards, not just what makes sense from experience.
\end{block}
\scenario{A bank identifies a foreign PEP customer with unexplained wealth. The bank has already performed EDD but cannot verify the source of wealth. The relationship manager wants to close the account.}
\vspace{0.2cm}
\textbf{Q: What is the bank's MOST appropriate action?}
\begin{itemize}
\item A) Close the account immediately
\item B) File a SAR and close the account
\item C) File a SAR, continue monitoring, and seek senior management approval to maintain or exit
\item D) Ignore the issue since EDD was already performed
\end{itemize}
\trap{Both B and C are partially correct (SAR filing is required). But immediate closure without approval may violate tipping-off.}
\correct{Correct Answer: C. Under FATF standards, exit should be risk-based and coordinated. Senior management must approve. SAR is mandatory regardless.}
\examfocus{Read for words like "most," "primary," "best," "first." These signal a "most correct" question.}
\end{frame}

% ============================================================
% SECTION 7: ADDITIONAL MISSED CONCEPTS
% ============================================================

\begin{frame}{Wolfsberg Group, Basel, Egmont - Specific Functions}
\begin{block}{Not Just "What They Are" - Specific Exam Angles}
Test-takers report questions on \textbf{specific secondary functions} of these bodies, not just primary missions.
\end{block}
\begin{tabular}{p{0.25\textwidth}|p{0.70\textwidth}}
\hline
\rowcolor{lightgray}
\textbf{Body} & \textbf{Exam-Tested Detail} \\
\hline
\textbf{Wolfsberg Group} & \textbf{NOT a regulator}. Association of 13 global banks. Produces \textbf{correspondent banking due diligence questionnaires} (CDDQs) and AML principles for private banking. \\
\hline
\textbf{Basel Committee} & Sets banking supervision standards (Basel III). Its \textbf{CDD paper} (1998) was foundational for KYC. Tests on history, not just capital rules. \\
\hline
\textbf{Egmont Group} & Network of FIUs. Facilitates \textbf{secure information exchange} between FIUs. NOT an investigative body. Exam tests: Egmont does NOT investigate; domestic FIUs do. \\
\hline
\end{tabular}
\redflag{QUESTION: "Which body produces the standard correspondent banking due diligence questionnaire?" → Wolfsberg Group.}
\end{frame}

\begin{frame}{Dual-Use Goods \& Proliferation Financing}
\begin{block}{Definition}
Dual-use goods are items that have legitimate civilian applications but can also be used in weapons of mass destruction (WMD) programs.
\end{block}
\textbf{Exam Red Flags:}
\begin{itemize}
\item Sale of dual-use goods to sanctioned countries (North Korea, Iran)
\item Transactions involving UNSC Resolution 1540 (non-proliferation)
\item Customer is a trading company with no apparent use for precision bearings, specialized valves, or aerospace-grade aluminum
\item Shipping routing through intermediary countries known for transshipment
\end{itemize}
\scenario{A small trading company orders \$2M in precision ball bearings rated for 10,000+ RPM. The buyer has never ordered industrial goods before. The destination is a front company in a country under UNSC sanctions for WMD programs.}
\vspace{0.1cm}
\textbf{Q: This is most indicative of?}
\correct{Correct Answer: Proliferation financing. The bank should file a SAR and block the transaction if sanctions apply.}
\end{frame}

\begin{frame}{Summary: Top 10 Missed Hard Questions Concepts}
\begin{enumerate}
\item \textbf{FATF IOs} (IO.1-IO.11): Know scenarios that map to each IO (especially IO.5, IO.6, IO.2)
\item \textbf{Life Insurance ML}: Policy loan = layering. CDD applies at \$10k premium or upon beneficiary designation
\item \textbf{Patriot Act Sections}: 311, 312, 313, 314(a), 314(b), 319(b), 326
\item \textbf{DeFi VASP Status}: >50\% governance tokens OR admin key = likely VASP
\item \textbf{Nested Correspondent}: Undisclosed nesting = SAR + RFI + potential termination
\item \textbf{"Most Correct" Questions}: Two correct answers; choose the one most aligned with FATF standards
\item \textbf{Wolfsberg Group}: Produces correspondent banking CDD questionnaires (not a regulator)
\item \textbf{Egmont Group}: Information exchange between FIUs; does NOT investigate
\item \textbf{Proliferation Financing}: Dual-use goods to sanctioned entities
\item \textbf{De-risking vs. Risk-Based Approach}: Blanket termination is NOT preferred; targeted exit is.
\end{enumerate}
\greennote{These concepts are consistently reported as "surprise" questions on recent CAMS exams. Master them to move from passing to confident.}
\end{frame}

\begin{frame}{Final Exam Day Reminders - From Recent Passers}
\begin{itemize}
\item \textbf{Time pressure is real}: 120 questions, 3.5 hours. Many questions require selecting 2-3 answers. Flag and move on.
\item \textbf{Version 6 vs. Version 7}: If taking V7, expect MORE on FATF IOs, less on pure memorization of recommendation numbers.
\item \textbf{Practice exams are EASIER than real exam}: Do not be overconfident from 90\% on ACAMS mock.
\item \textbf{Discord/ExamTopics}: Helpful but some answers are wrong. Verify against the study guide.
\item \textbf{Read the GLOSSARY}: Test-takers report the glossary saved them on obscure term questions.
\item \textbf{Insurance and Patriot Act}: Two areas that surprise experienced AML professionals.
\item \textbf{You can pass with 2-3 months of study while working full time.}
\end{itemize}
\vspace{0.5cm}
\centering
\greennote{Good luck - you've covered the hard parts now!}
\end{frame}

\end{document}